\chapter{PENUTUP}

Bab ini menjelaskan arah pengembangan penelitian di masa depan
berdasarkan keterbatasan penelitian yang telah dilakukan.

%%%%%%%%%%%%%%%%%%%%%%%%%%%%%%%%%%%%%%%%%%%%
\section{Keterbatasan Penelitian}
%%%%%%%%%%%%%%%%%%%%%%%%%%%%%%%%%%%%%%%%%%%%

Bagian ini menjelaskan keterbatasan yang terdapat dalam penelitian,
baik dari sisi dataset, metode, maupun lingkungan eksperimen.

Contoh keterbatasan penelitian:

\begin{itemize}

\item Dataset yang digunakan dalam penelitian masih terbatas pada
sumber tertentu sehingga belum mencakup seluruh kemungkinan
variasi data yang ada.

\item Pengujian sistem hanya dilakukan pada lingkungan eksperimen
tertentu sehingga perlu dilakukan evaluasi lebih lanjut pada
lingkungan yang lebih kompleks.

\end{itemize}

%%%%%%%%%%%%%%%%%%%%%%%%%%%%%%%%%%%%%%%%%%%%
\section{Arah Penelitian Selanjutnya}
%%%%%%%%%%%%%%%%%%%%%%%%%%%%%%%%%%%%%%%%%%%%

Bagian ini menjelaskan peluang penelitian lanjutan yang dapat dilakukan
berdasarkan hasil penelitian yang telah diperoleh.

Contoh arah penelitian selanjutnya:

\begin{enumerate}

\item Mengembangkan metode yang diusulkan dengan pendekatan algoritma
yang lebih kompleks.

\item Mengintegrasikan sistem dengan platform lain sehingga dapat
digunakan pada lingkungan yang lebih luas.

\item Melakukan evaluasi performa sistem menggunakan dataset dengan
skala yang lebih besar.

\end{enumerate}